\section{Discussion}\label{sec:discussion}

\paragraph{Why eleven rules nearly match logistic regression on
credit-card fraud.} The ULB features are PCA components; the fraud
signal concentrates in the tails of a handful of them, and quantile
propositions capture exactly those tails. Where the signal is
tail-shaped, a small rule set loses little discrimination and buys full
traceability. The synthetic study shows the opposite regime: when the
generating process is linear in the raw continuous features,
discretization pays a visible tax. The practical guidance for
audit-style deployments is therefore diagnostic, not dogmatic: measure
whether the domain's signal survives propositionalization (the static
rules' AUC against a raw-feature reference makes that a one-table
check) before committing to a rule-shaped model.

\paragraph{What maximum likelihood buys, and what it does not.} Across
every study, learning the weights improves calibration --- on the real
credit-card data the learned variant's log-loss is roughly a third below
the static rules' and its expected calibration error is smaller by a
factor of six. Ranking, by contrast, is already fixed by \emph{which
rules fire}: static noisy-OR and fitted weights order test cases almost
identically. In deployments where scores feed thresholds, queues, or
cost-sensitive decisions, calibration is the property that matters, and
it is the one learning delivers.

\paragraph{The intercept lesson generalizes.} The rank inversion of
Section~\ref{sec:learning} is not a quirk of one dataset: any additive
log-odds pooling of rules that are strong lifts but weak absolute
probabilities will invert on rare-event data unless a global intercept
absorbs the base rate. Static noisy-OR is immune because it is monotone
in the fired set --- more support can only raise the score. This
suggests a design rule for confidence calculi generally: monotonicity in
evidence is a rank-robustness property worth preserving, and learnable
calculi should carry an explicit prior term. The engine's export of the
intercept as a base-rate prior rule makes the repaired model exactly
representable in the same trace vocabulary.

\paragraph{What the drift studies say about context conditioning.} The
accounting experiments are the honest frontier of this paper, and their
verdict is negative. Under the strict pre/post-2003 split, rules mined
before the regime change barely rank later frauds; under the
SOX-boundary design --- built so the context weight is identifiable ---
the ablation gap is $0.4651$ vs.\ $0.464$: nothing. The mechanism
operates on rule \emph{reliability}, but what the regime change broke
was rule \emph{vocabulary}: antecedents that stopped being informative
cannot be rescued by reweighting, context-conditioned or otherwise.
When drift replaces the vocabulary, rule re-mining --- not reweighting
--- is the needed adaptation, connecting this line to concept-drift
taxonomies~\cite{gama2014} that distinguish real from virtual drift.
Whether richer context features (beyond the single-bit indicators
tested) change this verdict is the sharpest open question the negative
result leaves behind.

\paragraph{Traces as the deliverable.} The case study in
Section~\ref{sec:evaluation} shows what the auditor actually receives:
the fired rules with their empirical precisions, the aggregation
arithmetic line by line, and a deletion counterfactual quantifying the
top factor's contribution. None of the baselines --- including the
interpretable ones --- produce this artifact: a decision tree yields a
path, EBM yields shape functions, but neither names weighted domain
rules whose provenance is an auditable training-set measurement. The
fidelity results license trusting these traces for the static model;
for calibrated rare-event models the absolute-probability deletion
metrics compress, and log-odds-scale variants are the right refinement.
